\section{Implementation Details}
\label{app:impl}

\subsection{GNN Message Passing Primer}
\label{app:gnn-primer}
For completeness, we recall the message-passing form used by the GNN
backbones in this paper. Let $h_v^{(0)}=x_v$ be the feature vector of
node $v$. An $L$-layer GNN updates node representations by aggregating
messages from the graph neighborhood,
\begin{equation}
h_v^{(\ell+1)}
= \phi_\ell\!\left(h_v^{(\ell)},
  \operatorname{AGG}_{u\in\mathcal{N}(v)}
  \psi_\ell(h_v^{(\ell)},h_u^{(\ell)},e_{uv})\right),
\qquad \ell=0,\ldots,L-1,
\end{equation}
and predicts from $h_v^{(L)}$. Thus a deletion request is not an
i.i.d.\ sample removal: deleting a node removes its own training loss
and changes the message-passing context of retained neighbors. Exact
retraining can re-optimize under this modified graph, while approximate
unlearning must update from the original parameters. This is why our
diagnostics separate the effect of data removal from approximation drift
and report hop-distance decay relative to the GNN receptive field.

\subsection{Selector Details}
\label{app:selectors}

\begin{table}[h]
\centering
\scriptsize
\caption{\textbf{Four-tier access spectrum for selection strategies.}
Cost is measured after obtaining the deletion API. L2-\textsc{direct}
is an upper bound on the shadow-model setting.}
\label{tab:access}
\begin{tabular}{@{}lllc@{}}
\toprule
Tier & Information available & Cost & Strategies \\
\midrule
L0 \emph{(API only)}       & request $S\subset\mathcal{D}_{\mathrm{tr}}$        & $0$             & \textsc{Random} \\
L1 \emph{(structural)}     & public $G=(V,E)$, training mask                    & $0$ (coalition) & \textsc{Degree}, \textsc{PageRank}, \textsc{IM} \\
L2-\textsc{direct}         & production $f_\theta$, gradients, checkpoints      & internal leak   & \textsc{TracIn}, \textsc{Hybrid} \\
L2-\textsc{surrogate}      & shadow $\hat f_{\hat\theta}$ from public protocol  & GPU-h           & \textsc{TracIn}, \textsc{Hybrid} \\
L3 \emph{(white-box)}      & + modify training data / algorithm                 & internal access & \emph{out of scope} \\
\bottomrule
\end{tabular}
\end{table}

\paragraph{TracIn.}
Following~\cite{pruthi2020estimating}, we compute an influence proxy by
gradient inner-products between training and validation losses:
\begin{equation}
\label{eq:tracin}
\mathrm{TracIn}(v)=
\sum_c \eta_c
\left\langle
\nabla_\theta\mathcal{L}(f_{\theta_c};v),
\nabla_\theta\mathcal{L}(f_{\theta_c};V_{\mathrm{val}})
\right\rangle .
\end{equation}
The checkpoints $\theta_c$ are sampled periodically across training.  In this
revision, Eq.~\eqref{eq:tracin} denotes the paper-aligned trajectory selector
validated in the isolated gate of Appendix~\ref{app:approx-validity}.  The
historical main-matrix TracIn artifacts predate this contract and use a
deployed cross-final direction; they are not evidence for
Eq.~\eqref{eq:tracin} and require a versioned refresh before production use.

\paragraph{Influence Maximization.}
We treat the GNN receptive field as a diffusion graph and select nodes
maximizing reachable test disturbance under an independent-cascade model
with edge weight $p_{uv}=1/\deg(v)$. IM is model-free and therefore L1:
it depends on graph topology and degree, not on $f_\theta$.

\paragraph{Hybrid.}
Hybrid ranks candidates by
\begin{equation}
\label{eq:hybrid}
\mathrm{Score}_\alpha(v)
=\alpha\,\widetilde{\mathrm{TracIn}}(v)
 +(1-\alpha)\,\widetilde{\mathrm{IM}}(v),
\end{equation}
where $\widetilde{\cdot}$ denotes min-max normalization. We treat
$\alpha$ as a mixing coordinate: $\alpha=0$ is the IM endpoint,
$\alpha=1$ is the TracIn endpoint, and $\alpha=0.5$ is used in the main
matrix.  Because the historical Hybrid artifacts inherit the deployed
cross-final branch, any paper-aligned Hybrid conclusion remains gated on a
versioned trajectory-score parent artifact.

\paragraph{Batched CELF.}
We adopt CELF++~\cite{goyal2011celf++} with two modifications:
(a)~candidate-set pruning to the top-degree decile of the training set
(\texttt{candidate\_fraction=0.1}) before MC simulation, reducing the
search space by $\sim$10$\times$ on ogbn-arxiv with negligible spread loss;
(b)~a \emph{decoupled} Monte-Carlo seed
(\texttt{im\_selector\_seed}, default 2024) that is independent of the
GNN training seed, so that the selected node set is identical across the
five GNN seeds. The hyperparameter sweep is in
Appendix~\ref{app:imhp}; the selector-seed sanity test is in
Appendix~\ref{app:imseed}.

\subsection{TracIn for GNNs}
For node classification we treat the per-node training loss as
$\mathcal{L}_v = \mathrm{CE}(f_\theta(v), y_v)$ and the validation loss
as the mean over $V_{\mathrm{val}}$.  The isolated V2 gate uses five
post-update checkpoints (epochs 1, 10, 25, 50, and 100), weighted by the
learning rate of the update that produced each checkpoint.  This configuration
is prototype evidence, not a claim that the production runner already uses the
same artifact contract.

\subsection{Compute Resources}
Cora and Citeseer experiments fit on a single 4090 (24~GB); ogbn-arxiv
runs on an 80~GB H800 with batch size 256 and uses up to
${\sim}22$~GB peak GPU memory for the collateral retrain pass.
Total compute for the full Phase~B matrix (${\sim}1{,}500$ cells
across 4 configs) is approximately $14$--$30$~GPU-hours on the 4090
plus ${\sim}12$~GPU-hours on the H800.

\section{Additional Results}
\label{app:results}

This appendix collects per-cell tables, hop-decay curves per family,
and ogbn-arxiv feasibility commentary. The full per-seed aggregate
CSV ($N{=}460$ rows: the main matrix of 6 methods $\times$ 6
strategies $\times$ 5 seeds $\times$ \{Cora/GCN, Cora/GAT\} at
$r{=}0.05$, plus the 90 rows of the $r{=}0.01$ ablation
(Appendix~\ref{app:ratio}) and the 10 rows of the $\alpha{=}0$
endpoint (Appendix~\ref{app:alpha})) is shipped with the release as
\texttt{results/\_phase\_b\_aggregate.csv}. One cell---
GraphRevoker$\times$random$\times$seed42 at $r{=}0.01$---is flagged
\texttt{failed{=}true} (corrupted checkpoint, no valid
\texttt{f1\_after}); all paired effects in this appendix exclude it
from the random reference for that (method, seed) pair.

\subsection{1D Vulnerability Spectrum (Marginal Projection)}
\label{app:results-1d}
A 1D bar chart over methods using the $\ell_2$ norm of the
$(\Delta F^{\text{attack}}_{\mathrm{IM}},
\Delta F^{\text{attack}}_{\mathrm{TracIn}})$ fingerprint coordinate
preserves the prior ``vulnerability spectrum'' framing for warm-up
readers; the canonical ranking on Cora/GCN is GNNDelete $\gg$ GIF
$>$ MEGU $>$ IDEA $>$ GraphRevoker $\approx$ GraphEraser.

\subsection{Per-Cell Effect Tables}
\label{app:per-cell}
Tables~\ref{tab:per-cell-gcn} and~\ref{tab:per-cell-gat} report the
mean$\pm$std paired $\Delta F^{\mathrm{attack}}$ in F1 percentage
points across 5 seeds for Cora/GCN and Cora/GAT at $r{=}0.05$ for
all 6 methods $\times$ 6 strategies. The same numbers, with per-seed
values and the additional fields (\texttt{f1\_after\_mean},
\texttt{mia\_auc\_mean}, \texttt{gap\_mean\_pct}), are shipped as
\texttt{results/\_phase\_b\_aggregate.csv} (92 rows, including the
$r{=}0.01$ ablation cells of Appendix~\ref{app:ratio} and the
$\alpha{=}0$ endpoint of Appendix~\ref{app:alpha}). The ogbn-arxiv
3-seed feasible subset is added once the H800 retrain pass completes.

\begin{table}[h]
\centering
\footnotesize
\setlength{\tabcolsep}{4pt}
\caption{\textbf{Per-cell paired $\Delta F^{\mathrm{attack}}$ (\% pts,
mean$\pm$std over 5 seeds), Cora/GCN, $r{=}0.05$.} Random is the
reference (paired effect identically zero by construction).
\textbf{Bold} marks the strongest attack per row. Numbers reproduce
Table~\ref{tab:benchmark} with seed-level dispersion exposed; the
GraphRevoker row supersedes any earlier estimate (the dispatcher fix
of 2026-05-05 reassigned what was logged as ``GraphRevoker'' before
that date back to the GraphEraser branch).}
\label{tab:per-cell-gcn}
\begin{tabular}{@{}llrrrrr@{}}
\toprule
Family & Method & Degree & PageRank & IM & Hybrid & TracIn \\
\midrule
\multirow{2}{*}{Partition}
 & GraphEraser  & $\mathbf{+0.37}_{\pm 0.74}$ & $+0.44_{\pm 0.56}$ & $-0.00_{\pm 0.50}$ & $+0.29_{\pm 0.87}$ & $+0.04_{\pm 0.46}$ \\
 & GraphRevoker & $+0.11_{\pm 0.85}$ & $\mathbf{+0.33}_{\pm 0.84}$ & $+0.19_{\pm 0.89}$ & $+0.07_{\pm 0.87}$ & $-0.04_{\pm 1.09}$ \\
\midrule
\multirow{2}{*}{IF-based}
 & GIF  & $\mathbf{+2.14}_{\pm 0.21}$ & $+1.66_{\pm 0.34}$ & $+0.66_{\pm 0.43}$ & $+0.15_{\pm 0.63}$ & $-0.37_{\pm 0.60}$ \\
 & IDEA & $\mathbf{+1.26}_{\pm 0.46}$ & $+0.37_{\pm 0.49}$ & $+0.48_{\pm 0.67}$ & $+0.22_{\pm 0.53}$ & $-0.07_{\pm 0.55}$ \\
\midrule
\multirow{2}{*}{Learning}
 & GNNDelete & $\mathbf{+6.02}_{\pm 7.34}$ & $+5.20_{\pm 9.97}$ & $+3.21_{\pm 8.04}$ & $-1.95_{\pm 5.49}$ & $-4.61_{\pm 5.22}$ \\
 & MEGU      & $\mathbf{+1.74}_{\pm 0.28}$ & $+1.33_{\pm 0.20}$ & $+1.04_{\pm 0.21}$ & $+0.63_{\pm 0.46}$ & $-0.22_{\pm 0.66}$ \\
\bottomrule
\end{tabular}
\end{table}

\begin{table}[h]
\centering
\footnotesize
\setlength{\tabcolsep}{4pt}
\caption{\textbf{Per-cell paired $\Delta F^{\mathrm{attack}}$ (\% pts,
mean$\pm$std over 5 seeds), Cora/GAT, $r{=}0.05$.} Same conventions
as Table~\ref{tab:per-cell-gcn}.}
\label{tab:per-cell-gat}
\begin{tabular}{@{}llrrrrr@{}}
\toprule
Family & Method & Degree & PageRank & IM & Hybrid & TracIn \\
\midrule
\multirow{2}{*}{Partition}
 & GraphEraser  & $+0.15_{\pm 1.59}$ & $\mathbf{+0.89}_{\pm 2.36}$ & $+0.63_{\pm 1.11}$ & $-0.22_{\pm 1.63}$ & $+0.30_{\pm 2.11}$ \\
 & GraphRevoker & $\mathbf{+1.44}_{\pm 1.23}$ & $+1.29_{\pm 1.05}$ & $+1.25_{\pm 1.45}$ & $+1.37_{\pm 1.74}$ & $+0.22_{\pm 1.99}$ \\
\midrule
\multirow{2}{*}{IF-based}
 & GIF  & $\mathbf{+2.88}_{\pm 2.12}$ & $+2.22_{\pm 1.99}$ & $+1.11_{\pm 1.34}$ & $+1.77_{\pm 1.87}$ & $+0.78_{\pm 1.62}$ \\
 & IDEA & $\mathbf{+3.32}_{\pm 0.84}$ & $+2.69_{\pm 1.31}$ & $+1.11_{\pm 0.83}$ & $+0.81_{\pm 1.06}$ & $-0.11_{\pm 1.41}$ \\
\midrule
\multirow{2}{*}{Learning}
 & GNNDelete & $\mathbf{+1.73}_{\pm 5.78}$ & $+1.18_{\pm 3.36}$ & $-0.07_{\pm 3.28}$ & $-0.74_{\pm 2.13}$ & $-0.00_{\pm 2.10}$ \\
 & MEGU      & $\mathbf{+1.10}_{\pm 0.72}$ & $+0.70_{\pm 0.72}$ & $+0.07_{\pm 0.72}$ & $+0.11_{\pm 0.89}$ & $-0.45_{\pm 0.84}$ \\
\bottomrule
\end{tabular}
\end{table}

\paragraph{Reading.} Three patterns are visible at the per-cell
resolution that the headline benchmark cannot show: (i)~the
\textsc{Degree} column is the strongest attack in 11 of 12 rows
(the lone exception is GraphEraser/GAT, where \textsc{PageRank}
narrowly leads); the L1-structural baseline is consistently the
upper-envelope selector. (ii)~The \textsc{TracIn} column is
non-positive in 8 of 12 rows---the L2-\textsc{direct} gradient
selector is structurally weaker than degree, in line with the
selection-degree mechanism of \S\ref{sec:results-alignment}.
(iii)~Variance is concentrated in two cells: GNNDelete on both
backbones (std~5--10\%, the unstable-collapse signature) and
GraphRevoker/GAT$\times$\textsc{TracIn} (std~2.0\%, the lone
significant non-degree cell of \S\ref{sec:results-fingerprint}).
Other rows are tight (std~$<1\%$).

\subsection{Hop-Distance Decay Curves}
On Cora/GCN ($L{=}2$) under TracIn-based deletion, GIF and IDEA
flips are sharply concentrated within the 1-hop receptive field
(${\approx}7\%$ flip rate at $h{=}1$, ${<}0.3\%$ at $h{\ge}2$);
GNNDelete is far more diffuse (${\approx}48\%$ at $h{=}1$,
${\approx}18\%$ at $h{=}2$, $4.5\%$ at $h{=}3$). The Partition pair
shows broad flips across all hop buckets ($35$--$40\%$ at $h{=}1$,
$20\%$+ even at $h{>}3$), consistent with shard re-training rather
than localised approximation drift.

\section{Influence-Selector Taxonomy and Concordance}
\label{app:if-gif-taxonomy}

We separate the influence-family selectors by the object they score.
The target-free A/B group measures gradient magnitude
$\lVert g_v\rVert$ or curvature-adjusted parameter displacement
$\lVert H^{-1}g_v\rVert$.  The C group contains evaluation-conditioned
IF scores whose source is either $g_v$ or the original-graph affected-set
gradient $\mathrm{grad1}_v$.  The D group is reserved for graph-aware GIF
and therefore uses the full deletion source
$\mathrm{grad1}_v-\mathrm{grad2}_v$.  Legacy cross-TracIn and topology
selectors are external controls rather than D-group methods.

\begin{table}[h]
\centering
\small
\caption{\textbf{Within- and cross-group selector concordance.}
Spearman correlations are descriptive means over three Planetoid datasets
and three seeds with a two-layer GCN.  A/B, C, and D denote parameter-change,
IF, and graph-aware GIF groups, respectively.}
\label{tab:if-gif-concordance}
\begin{tabular}{@{}llr@{}}
\toprule
Test & Comparison & Spearman \\
\midrule
A/B proxy & A vs. B reference & $0.962$ \\
C/D separation & C-simple vs. D-full & $0.040$ \\
D final proxy & P-graph vs. D-full & $0.984$ \\
D trajectory & graph CP-3 / CP-6 vs. D-full & $0.498/0.529$ \\
\bottomrule
\end{tabular}
\end{table}

The A/B correlation shows that gradient magnitude is a strong
Hessian-free proxy for the B ranking in this setting; it does not make
the two definitions equivalent.  All inverse-Hessian terms are computed
through iterative IHVP solves; solver, HVP, and probe choices are
reproducibility details rather than separate selector families.  The
near-zero C-simple/D-full correlation shows that adding
the post-deletion graph term is a mechanism change, whereas the high
P-graph/D-full correlation shows that the Hessian can be removed reliably
when the D source is held fixed.  B reference versus D-full ($0.023$) and
C-point versus D-full ($0.112$) provide additional cross-group separation.

This study uses $3$ datasets $\times$ $3$ seeds and deletion budgets
$k\in\{3,7,14\}$.  It evaluates local rankings and set-deletion outcomes;
it is not an end-to-end graph-unlearning comparison and supports no
significance claim.  Production registration of the proper-TracIn selector
is governed by a separate cache/recipe gate.

\section{Approximation Validity of IM and Trajectory Influence Proxies}
\label{app:approx-validity}

We next ask whether two proxy objectives preserve the signals they claim to
represent.  This is distinct from Appendix~\ref{app:if-gif-taxonomy}: the
previous section separates influence targets, whereas this section validates
representative approximations within their intended axes.

\paragraph{IM coverage is not degree in disguise.}
Table~\ref{tab:im-degree-validity} reports top-set overlap between IM and
degree on six datasets.  IM has low Jaccard overlap on every dataset and is
most distinct on the larger or denser graphs.  Thus its coalition-level
coverage signal does not reduce to a pointwise degree ranking, although this
does not imply stronger downstream attack damage.

\begin{table}[h]
\centering
\scriptsize
\caption{\textbf{Set-level validity of the IM coverage proxy.} Jaccard
overlap between the top IM and degree selections at the recorded 5\% budget.}
\label{tab:im-degree-validity}
\begin{tabular}{@{}lrrrrrrr@{}}
\toprule
Dataset & Cora & CiteSeer & PubMed & Photo & Computers & CS & Mean \\
\midrule
IM vs. degree & 0.187 & 0.177 & 0.061 & 0.046 & 0.030 & 0.032 & 0.089 \\
\bottomrule
\end{tabular}
\end{table}

\paragraph{Trajectory influence is aligned but not uniformly so.}
The isolated V2 gate compares the multi-checkpoint score in
Eq.~\eqref{eq:tracin} with an evaluation-conditioned IF reference over
$3$ Planetoid datasets, $2$ backbones, and $3$ seeds.
Table~\ref{tab:tracin-v2-validity} shows that configuration-mean Spearman
correlations range from $0.762$ to $0.962$.  PubMed/GAT is the explicit
sensitivity case: its common Top-7 counts fall to $[4,4,1]$ across seeds, so
the evidence supports typical alignment rather than universal equivalence.

\begin{table}[h]
\centering
\small
\caption{\textbf{Trajectory-proxy alignment with eval-IF.} Spearman is the
three-seed configuration mean; brackets report common Top-7 nodes for seeds
$[2024,7,42]$.}
\label{tab:tracin-v2-validity}
\begin{tabular}{@{}lrr@{}}
\toprule
Configuration & Spearman & Common@7 \\
\midrule
Cora / GCN     & 0.877 & $[5,6,5]$ \\
CiteSeer / GCN & 0.950 & $[6,5,4]$ \\
PubMed / GCN   & 0.908 & $[5,3,5]$ \\
Cora / GAT     & 0.825 & $[6,7,5]$ \\
CiteSeer / GAT & 0.962 & $[6,6,5]$ \\
PubMed / GAT   & 0.762 & $[4,4,1]$ \\
\bottomrule
\end{tabular}
\end{table}

\paragraph{Evidence boundary.}
The IM result is a deterministic set-level selector comparison, and the
trajectory result is an isolated conditional-prototype gate.  Neither measures
graph-unlearning damage.  Formal immutable score storage, Legacy-store
isolation, Hybrid parent identity, and an end-to-end GU canary remain required
before the trajectory selector can support production or main-matrix claims.

\section{Ablation: IM Selector Seed Decoupling}
\label{app:imseed}

\textbf{Purpose.} Validate that the Vulnerability Fingerprint and
$\alpha$-sweep findings are not artefacts of selector noise.
Table~\ref{tab:jaccard} reports selection-set Jaccard before and after
the seed-decoupling fix.

\begin{table}[h]
\centering
\small
\caption{\textbf{Selection-set Jaccard across 5 GNN seeds, before vs.\
after the IM selector-seed decoupling fix.} Cora/GCN/$r{=}0.05$.}
\label{tab:jaccard}
\begin{tabular}{@{}lrr@{}}
\toprule
Strategy & Pre-fix (coupled) & Post-fix (decoupled) \\
\midrule
\textsc{Random}    & $0.023$ & $0.023$ \\
\textsc{Degree}    & $1.000$ & $1.000$ \\
\textsc{PageRank}  & $1.000$ & $1.000$ \\
\textsc{TracIn}    & $0.415$ & $0.415$ \\
\textsc{IM}        & $0.129$ & $\mathbf{1.000}$ \\
\textsc{Hybrid}    & $0.646$ & $0.917$ \\
\bottomrule
\end{tabular}
\end{table}

After the fix, the paired-effect 95\% CI width on the GNNDelete IM
cell is $\pm 9.99\%$ ($N{=}5$, two-sided $t$-CI), reflecting the
inherently high seed-to-seed variance of GNNDelete's collapse mode
rather than selector-induced jitter.

\section{Ablation: GraphEraser Shard Configuration}
\label{app:shard}

\textbf{Purpose.} Lock down Shard Protection
(\S\ref{sec:results-shard}) as a mechanism-driven property rather than
a hyperparameter coincidence. We sweep shard counts
$C\in\{5,10,20\}$ on Cora/GCN/$r{=}0.05$ with strategies
$\{\textsc{Random},\textsc{IM},\textsc{Hybrid}\}$ and 5 seeds. The
paired effect remains ${\le}0$ across all shard counts; the
noise-floor term $\Delta F_{\text{noise}}$ stays negative
($-5\%$ to $-15\%$ on Cora), confirming Shard Protection is robust to
the partition-count hyperparameter.

\section{Ablation: \texorpdfstring{$\alpha$}{alpha}-Sweep as Synergy Diagnostic}
\label{app:alpha}

\textbf{Purpose.} We treat $f(\alpha)\!:=\!\Delta F^{\text{attack}}$ of
the Hybrid score (Eq.~\ref{eq:hybrid}) as a synergy diagnostic; the
curve shape encodes how IM and TracIn signals interact per family.
Curve readings: linear $\Rightarrow$ independent additive selection;
concave (above chord) $\Rightarrow$ \emph{synergy}; convex (below
chord) $\Rightarrow$ \emph{redundancy}; flat at zero $\Rightarrow$
mechanism-incomparable.

\textbf{Design.} 5 $\alpha$ values $\{0.0, 0.25, 0.5, 0.75, 1.0\}$,
6 methods, 2 cells (Cora/GCN, Cora/GAT), 5 seeds; the endpoints
$\alpha\in\{0,1\}$ are reused from the main matrix. Total new runs
$N{=}180$. The full $\alpha$-sweep grid is deferred to a
supplementary release; the $\alpha{=}0.5$ column appears in the main
fingerprint as the ``hybrid'' strategy.

\paragraph{$\alpha{=}0$ endpoint (IF-only Hybrid path).} A
spot-check ablation runs the Hybrid pipeline with the IM weight
zeroed out ($\alpha{=}0$), so the fused score reduces to a
TracIn-only ranking under the Hybrid rescaling. On Cora/GCN/$r{=}0.05$
the paired effect at $\alpha{=}0$ is $+1.07_{\pm 0.68}\%$ for GIF and
$+5.46_{\pm 5.84}\%$ for GNNDelete (5 seeds, paired vs.\ same-seed
random). For comparison, the plain \textsc{TracIn} strategy on the
same cells reads $-0.37_{\pm 0.60}\%$ (GIF) and $-4.61_{\pm 5.22}\%$
(GNNDelete). The $\sim 10\%$ pt swing on GNNDelete between the two
``IF-only'' selectors is mechanically traceable to the Hybrid path's
score normalisation (which rescales the TracIn score against the IM
score distribution before weighting), and it confirms that the curve
endpoint at $\alpha{=}0$ is \emph{not} interchangeable with the
plain \textsc{TracIn} column of Table~\ref{tab:per-cell-gcn}. The
$\alpha{=}0$ result is reported here for transparency; we use the
plain \textsc{TracIn} numbers in the main text and treat the Hybrid
$\alpha$-curve as a relative synergy diagnostic, not as an absolute
TracIn estimate.

\section{Ablation: IM Hyperparameters}
\label{app:imhp}

\textbf{Purpose.} Defend the ``${\sim}10\times$ speedup with
$<$2\% spread loss'' claim of \S\ref{sec:experiment} and the IM
scaling story of \S\ref{sec:setup-arxiv}. We sweep candidate-fraction
$\in\{0.05,0.1,0.2\}$ and MC rounds $\in\{25,50,100\}$ on
Cora/GCN/GIF/$r{=}0.05$, reporting paired effect, selection-set
Jaccard with full-CELF reference, and wall-clock. The full robustness
grid is shipped with the supplement; the canonical setting
(\texttt{candidate\_fraction}${=}0.1$, $50$ MC rounds) maintains
$\mathrm{Jaccard}{>}0.95$ against full-CELF on the same seed while
incurring $<2\%$ spread loss.

\section{Ablation: Deletion Ratio Sweep}
\label{app:ratio}

\textbf{Purpose.} Characterise per-family budget elasticity, removing
the ``$r{=}0.05$ only'' restriction.

\textbf{Design.} $r\in\{0.01,0.05\}$ (Phase~B.1; the
$r\in\{0.10,0.20\}$ extension is deferred to the supplementary
release), all 6 methods, strategies
$\{\textsc{Random}, \textsc{IM}, \textsc{TracIn}\}$, Cora/GCN,
5 seeds. The $r{=}0.01$ slice is $N{=}90$ runs.

\begin{table}[h]
\centering
\footnotesize
\setlength{\tabcolsep}{5pt}
\caption{\textbf{Paired $\Delta F^{\mathrm{attack}}$ (\% pts,
mean$\pm$std over 5 seeds) at $r{=}0.01$ vs $r{=}0.05$, Cora/GCN.}
For non-GNNDelete methods the $r{=}0.01$ effect is sub-1\% across
all selectors---the deletion budget is too small to perturb the
unlearning operator. GNNDelete is the exception: \textsc{IM} at
$r{=}0.01$ delivers a paired $+10.1\%$, three times its $r{=}0.05$
value. The ratio sweep therefore inverts the headline ranking on
this single cell---\textsc{IM} dominates \textsc{TracIn} only when
the volume signal of \S\ref{sec:results-shard} is suppressed.
GraphRevoker/seed42 was excluded by the gate due to a corrupted
checkpoint; values reflect the 4 valid seeds for the random column
of GraphRevoker only, all other cells are 5-seed.}
\label{tab:ratio-sweep}
\begin{tabular}{@{}llrrrr@{}}
\toprule
& & \multicolumn{2}{c}{$r{=}0.01$} & \multicolumn{2}{c}{$r{=}0.05$} \\
\cmidrule(lr){3-4}\cmidrule(lr){5-6}
Family & Method & IM & TracIn & IM & TracIn \\
\midrule
\multirow{2}{*}{Partition}
 & GraphEraser  & $+0.19_{\pm 0.58}$ & $-0.11_{\pm 0.48}$ & $-0.00_{\pm 0.50}$ & $+0.04_{\pm 0.46}$ \\
 & GraphRevoker & $-2.44_{\pm 6.40}$ & $-2.40_{\pm 6.50}$ & $+0.19_{\pm 0.89}$ & $-0.04_{\pm 1.09}$ \\
\midrule
\multirow{2}{*}{IF-based}
 & GIF  & $-0.00_{\pm 0.29}$ & $-0.37_{\pm 0.29}$ & $+0.66_{\pm 0.43}$ & $-0.37_{\pm 0.60}$ \\
 & IDEA & $+0.52_{\pm 0.36}$ & $-0.11_{\pm 0.21}$ & $+0.48_{\pm 0.67}$ & $-0.07_{\pm 0.55}$ \\
\midrule
\multirow{2}{*}{Learning}
 & GNNDelete & $\mathbf{+10.11}_{\pm 4.13}$ & $-0.85_{\pm 2.81}$ & $+3.21_{\pm 8.04}$ & $-4.61_{\pm 5.22}$ \\
 & MEGU      & $+0.37_{\pm 0.26}$ & $+0.04_{\pm 0.16}$ & $+1.04_{\pm 0.21}$ & $-0.22_{\pm 0.66}$ \\
\bottomrule
\end{tabular}
\end{table}

\textbf{Reading.} The $r{=}0.01$ slice cleanly bisects the Phase~B
methods into two mechanism classes.

\textit{Volume-driven (five of six methods).} GIF, IDEA, MEGU,
GraphEraser, and GraphRevoker all post paired
$|\Delta F^{\mathrm{attack}}|<1\%$ at $r{=}0.01$ on both \textsc{IM}
and \textsc{TracIn}, with std also $<1\%$ for the non-Partition
members. This is a stronger statement than the $k{=}5$ noise floor
of \S\ref{sec:results-shard}: it confirms that even sophisticated
selectors fail to extract additional damage at low deletion budget.
Combined with the substantial $r{=}0.05$ effects on the same
methods (Table~\ref{tab:per-cell-gcn}), this shows that for the
volume-driven family the attack-budget elasticity is super-linear,
and the apparent ``vulnerability'' at $r{=}0.05$ is mostly
volume-driven---exactly what
the same-seed paired baseline of \S\ref{sec:results-shard} already
indicates. GraphRevoker is the only Partition member whose
$r{=}0.01$ paired effect is sizeable in magnitude
($-2.4_{\pm 6.4}\%$), but the std swallows the mean and the sign is
\emph{negative} (Shard Protection holds at low budget too).

\textit{Architectural (GNNDelete only).} GNNDelete is the lone
exception. Its paired \textsc{IM} effect at $r{=}0.01$ is
$+10.11_{\pm 4.13}\%$---larger than its $r{=}0.05$ \textsc{IM} cell
($+3.21_{\pm 8.04}\%$) and the cleanest GNNDelete attack signal in
the entire Phase~B matrix (std cut roughly in half). Equally
striking: the $r{=}0.01$ \textsc{TracIn} cell is
$-0.85_{\pm 2.81}\%$, so the headline ``\textsc{TracIn} loses to
\textsc{Degree}/\textsc{IM}'' ordering of \S\ref{sec:results-fingerprint}
is preserved here even though every other method has flatlined.
Because the same-seed random baseline at $r{=}0.01$ delivers a near-zero
F1 drop on GNNDelete (volume effect is small at $14$ deletions out
of $\sim 2700$), the IM result is not a residual of the volume
signal---it is GNNDelete's mask network failing on a
selector-specific structural target.

\textit{Combined reading.} The $r{=}0.01$ ablation thus delivers a
two-class taxonomy of unlearning fragility that is not visible at
$r{=}0.05$ alone: GNNDelete is structurally vulnerable, the other
five methods only appear vulnerable insofar as the deletion budget
itself is large enough to perturb the model. The
$r\in\{0.10,0.20\}$ extension is deferred to the supplementary
release; its only open question is whether GNNDelete's collapse
continues to grow super-linearly past $r{=}0.05$ or saturates.

\section{Per-Method Hyperparameters}
\label{app:hyperparams}
We adopt the published defaults from each method's original paper
where specified, and our reproduced defaults (matching the OpenGU
benchmark) where not. The full per-method hyperparameter table is
included with the released configuration files
(\texttt{model/properties/}).

\section{Negative and Boundary-Case Analyses}
\label{app:negative}

\paragraph{Cases where attacks underperform random.}
On the partition pair (GraphEraser, GraphRevoker), all model-coupled
selectors yield smaller F1 gain than random---the Shard Protection
upper-bound (\S\ref{sec:results-shard}). The mechanism: adversarial
nodes concentrate in fewer shards than random selection, leading to
weaker aggregator disturbance; per-shard diagnostics confirming this
concentration are included in the supplementary release.

\section{Discussion}
\label{sec:conclusion}

\subsection{Interpreting Null Effects}
\label{sec:disc-outliers}

MEGU and IDEA should not be read as a new ``robust family'': they belong
to different OpenGU categories. Their shared property is geometric and
diagnostic---both decouple from the IM/TracIn selector signals tested
here. This is partly evidence of method-specific robustness and partly a
gap in the current attacker class. Mechanism-aware selectors targeting
mutual-evolution or certified-gradient objectives are the natural next
test.

\subsection{Discussion: Lessons for Secure GU Design}
\label{sec:discussion-design}

The observed family patterns offer actionable lessons for designing next-generation approximate GU methods that are both fast and secure. 
First, \textbf{Structural Isolation:} The Shard Protection effect shows that physical partitioning provides a natural firewall against cascading failure. Future learning-based or IF-based methods could adopt hybrid architectures where gradient updates are confined by soft structural boundaries. 
Second, \textbf{Gradient Regularization:} Methods sensitive on the TracIn axis (like GIF) highlight the risk of relying purely on unconstrained gradient similarity. Incorporating adversarial regularization during the approximation step—essentially making the unlearning step minimax robust to the worst-case deletion batch—could flatten this vulnerability. 
Ultimately, our fingerprint shows that efficiency (speedup) and utility retention under random deletion no longer suffice as the sole evaluation criteria; worst-case auditing must become standard practice.

\subsection{ogbn-arxiv Pilot Scale Check}
\label{app:arxiv-pilot}

Table~\ref{tab:arxiv-pilot} reports the late-arriving ogbn-arxiv pilot
subset used as a scale sanity check. Because this run covers only seed
42, two methods, and three selectors, we keep it outside the main
benchmark ranking and use it only to check whether the small-graph
fingerprint remains qualitatively plausible at 169K nodes. The pilot is
consistent with the main Cora/Citeseer reading: GIF is nearly invariant
under the tested selectors, while GNNDelete exhibits large utility loss
and retrain-gap drift, with IM adding an $18.17$ percentage-point excess
drop over budget-matched random deletion.

\begin{table}[t]
\centering
\footnotesize
\setlength{\tabcolsep}{5pt}
\caption{\textbf{ogbn-arxiv pilot subset.}
GCN backbone, $r{=}0.01$, seed 42 only. This late-stage scale sanity
check is not included in the main benchmark ranking because it covers
only two methods and three selectors. Total drop is
$F_1^{\mathrm{before}}-F_1^{\mathrm{unlearn}}$ in percentage points;
paired effect is measured relative to the same-method random baseline.}
\label{tab:arxiv-pilot}
\begin{tabular}{llrrrr}
\toprule
Method & Strategy & Total drop & Paired vs. random & Gap & Flip rate \\
       &          & (pp)       & (pp)              & (\%) & (\%) \\
\midrule
GIF       & Random & $-0.09$ & $0.00$   & $0.25$ & $4.82$ \\
GIF       & IM     & $-0.06$ & $+0.03$  & $0.35$ & $5.14$ \\
GIF       & TracIn & $-0.02$ & $+0.07$  & $0.14$ & $4.99$ \\
\midrule
GNNDelete & Random & $32.27$ & $0.00$   & $46.23$ & $53.72$ \\
GNNDelete & IM     & $50.44$ & $+18.17$ & $72.22$ & $75.52$ \\
GNNDelete & TracIn & $39.31$ & $+7.04$  & $56.25$ & $62.12$ \\
\bottomrule
\end{tabular}
\end{table}

\subsection{Limitations}
\label{sec:limitations}

\textbf{L1.} Shard-count coverage is partial (\S\ref{app:shard}); a
full partition-design sweep is future work.
\textbf{L2.} ogbn-arxiv uses 3 seeds rather than 5; we report CI widths
to make the power tradeoff explicit.
\textbf{L3.} TracIn's $G$-matrix does not scale to very high feature
dimension (Physics is out of scope).
\textbf{L4.} Our L2 axis reports L2-\textsc{direct} as an upper bound on
the L2-\textsc{surrogate} variant available to a coalition that trains
its own shadow model. Standard adversarial-transferability
work~\cite{papernot2017practical,demontis2019adversarial} suggests this
upper bound is rigorous in the worst case, but cross-architecture
transfer on graphs is not quantified here and is left to future work
(see also \S\ref{sec:future}).
\textbf{L5.} Ratio sweeps are small-graph only
($r\in\{0.01,0.05,0.10,0.20\}$ on Cora/GCN).
\textbf{L6.} We cover three of four OpenGU categories; ``Others'' and
third members within each family are deferred.
\textbf{L7.} The design spans seven axes, so we use claim-conditioned
slices rather than a Cartesian-complete grid; untested combinations are
reported as out of scope in Appendix~\ref{app:results}.
\textbf{L8.} The reported update-detection AUC is a posterior-shift
deletion-membership audit, not a full shadow-model MIA; scaling standard
shadow-model MIA across all GU families is left to future work.

\subsection{Broader Impact}
\label{sec:broader}

Approximate GU papers should report not only random-deletion utility and
runtime, but also adversarial-deletion bounds. The diagnostic kit here
($\Delta F$ decomposition, retrain gap, prediction shift, hop-distance
decay, and update-detection AUC) can be reused with stronger selectors.

A coalition-of-users can reach the L1 structural tier at zero additional
model-access cost; L2-\textsc{direct} should be read as an upper bound on
what a shadow-model attacker may realize after transfer. Shard Protection
is therefore a defensive insight rather than a deployment guarantee: it
limits utility collapse, but it does not by itself address privacy
objectives.

\subsection{Future Directions}
\label{sec:future}
Four problems remain: mechanism-aware selectors for MEGU/IDEA-like
objectives; defenses that quantify the privacy--robustness tradeoff of
partition-based GU (a natural baseline being a deletion-budget cap that
falls back to exact retraining once a coalition exceeds threshold);
cross-architecture surrogate transferability of model-coupled selectors,
which would tighten the L2-\textsc{direct}/\textsc{surrogate} gap of L4;
and attacks whose objective is privacy leakage or deletion visibility
rather than downstream F1.
